\documentclass{exam}
\usepackage{../commonheader}

\discnumber{5}
\title{Recursion}
\date{February 17 -- February 21, 2026}


\begin{document}
\maketitle

\begin{meta}
    \textbf{Example Timeline}
    \begin{itemize}
        \item Recursion Basics [4 min]
        \subitem Recursion, imo, requires more practice to get the jist of, so I would recommend spending more time walking through the first two problems as opposed to a longer mini-lecture.
        \subitem Make sure mentees understand the basic template of a basic case + recursive case 
        \subitem Some common metaphors/analogies for this: dividing a pizza, factorials, etc.
        \item 1. Num-Digits [6 min]
        \item 2. Right-Triangle [7 min]
        \item 3. Replace [10 min]
        \item Tree Recursion [5 min]
        \subitem introduce the idea of a recursive tree, what does depth mean? 
        \item 4. Gibonacci [8 min] 
        \item 5. Discussions1 [7 min]
        \item 6. Positioner [15 min] - optional, only attempt for advanced students
    \end{itemize}
\end{meta} 

\section{Recursion}
\begin{questions}
    \subimport{../../topics/recursion/easy/}{num-digits.tex}
    \begin{questionmeta}
        Suggested Time: 6 min; Difficulty: Easy
        \begin{itemize}
            \item  Good fundemental question to go over recursion basic template, intuition behind recursive call
        \end{itemize}
      \end{questionmeta}    
      
      \subimport{../../topics/recursion/easy/}{right_triangle.tex}
      \begin{questionmeta}
          Suggested Time: 4 min; Easy
          \begin{itemize}
              \item  Nice way to visualize recursion and introduces students to the idea that recursion may not always involve returning
              \item HIGHLY HIGHLY recommend asking follow up questions about what would happen if we changed parts of the code, or if we wanted to output a reversed triangle instead
          \end{itemize}
        \end{questionmeta}

        \subimport{../../topics/recursion/medium/}{replace.tex}
        \begin{questionmeta}
            Suggested Time: 9 min; Medium
            \begin{itemize}
              \item  Think about how to process numbers digit by digit. Explain how the *10 rebuilds the number digit by digit as the recursion unwinds.
          \end{itemize}
        \end{questionmeta}
\end{questions}

\section{Tree Recursion}
\begin{questions}
    \subimport{../../topics/recursion/easy/}{gibonacci.tex}
    \begin{questionmeta}
        Suggested Time: 9 min; Difficulty: Easy/Medium
        \begin{itemize}
            \item  Good introduction problem to tree recursion, especially with the similarity with fibonacci. 
            \item Make sure students understand how the base case works. 
        \end{itemize}
      \end{questionmeta}    
      
      \subimport{../../topics/recursion/medium/}{discussions1.tex}
      \begin{questionmeta}
          Suggested Time: 9 min; Medium
        \end{questionmeta}

        \subimport{../../topics/recursion/medium/}{positionizer.tex}
        \begin{questionmeta}
            Suggested Time: 15 min; Hard
          \end{questionmeta}


        
\end{questions}


\end{document}

